
\documentclass[11pt]{article}

\usepackage[margin=1in]{geometry}
\usepackage{amsmath,amssymb}
\usepackage{listings}
\usepackage{xcolor}
\usepackage{tikz}
\usepackage{hyperref}

\definecolor{backcolour}{rgb}{0.95,0.95,0.92}

\lstset{
    backgroundcolor=\color{backcolour},
    basicstyle=\ttfamily\small,
    breaklines=true,
    frame=single
}

\title{CPSC 250 \\ Functions}
\author{Edward Brash}
\date{}

\begin{document}

\maketitle

\section{Introduction}

Functions are one of the most important organizational tools in programming.

Functions allow us to:

\begin{itemize}
    \item divide problems into smaller pieces,
    \item reuse code,
    \item improve readability,
    \item and reduce duplication.
\end{itemize}

\section{Basic Function Structure}

General syntax:

\begin{lstlisting}[language=Python]
def my_function(arguments):

    # code here

    return values
\end{lstlisting}

Important parts:

\begin{itemize}
    \item Function name
    \item Argument list
    \item Function body
    \item Return value list
\end{itemize}

\subsection{Argument List}

The argument list contains information passed \emph{to} the function.

\subsection{Return Value List}

The return value list contains information returned \emph{from} the function.

\section{Conceptual Model}

Conceptually:

\begin{center}
Main Program $\rightarrow$ Function

\vspace{0.3cm}

Arguments flow into the function.

\vspace{0.3cm}

Return values flow back out.
\end{center}

\section{Example: add() Function}

\begin{lstlisting}[language=Python]
def add(a, b):

    c = a + b

    return c
\end{lstlisting}

Example usage:

\begin{lstlisting}[language=Python]
x = 4
y = 6

z = add(x, y)

print(z)
\end{lstlisting}

Output:

\begin{lstlisting}
10
\end{lstlisting}

\section{Very Important Idea: Pass by Value}

Inside the function:

\begin{itemize}
    \item \texttt{a} and \texttt{b} are copies of \texttt{x} and \texttt{y}
    \item they are different variables in memory
\end{itemize}

Changing variables inside the function does not change variables in the main program.

This is called:

\[
\textbf{pass by value}
\]

\section{Naming Convention Recommendation}

A very good habit is to use different variable names inside functions than in the main program.

Example:

\begin{center}
\begin{tabular}{|c|c|}
\hline
Main Program & Function \\
\hline
x, y, z & a, b, c \\
\hline
\end{tabular}
\end{center}

This reinforces the idea that these are different variables occupying different memory locations.

\section{Dynamic Typing in Python}

Unlike languages such as:

\begin{itemize}
    \item C
    \item C++
    \item Java
    \item Go
\end{itemize}

Python uses dynamic typing (sometimes informally called ``duck typing'').

The same function may operate on multiple data types.

Example:

\begin{lstlisting}[language=Python]
def add(a, b):

    return a + b
\end{lstlisting}

This works for integers:

\begin{lstlisting}[language=Python]
print(add(4, 6))
\end{lstlisting}

but also for strings:

\begin{lstlisting}[language=Python]
x = "that is "
y = "so cool!"

z = add(x, y)

print(z)
\end{lstlisting}

Output:

\begin{lstlisting}
that is so cool!
\end{lstlisting}

Python automatically determines the correct behavior based on object types.

\section{Variable Scope}

Scope determines where variables have meaning inside a program.

\subsection{Example}

\begin{lstlisting}[language=Python]
employee_name = 'Bob'

def get_name():

    name = input()

    employee_name = name

get_name()

print(employee_name)
\end{lstlisting}

Important:

The variable:

\begin{lstlisting}[language=Python]
employee_name
\end{lstlisting}

inside the function is different from the variable outside the function.

Thus, the printed value remains:

\begin{lstlisting}
Bob
\end{lstlisting}

unless special mechanisms such as \verb|global| are used.

\section{Pass by Reference?}

Students often hear the phrase:

\[
\textit{``Python passes lists by reference.''}
\]

This is not quite correct.

\section{The Common Misconception}

Suppose we have:

\begin{lstlisting}[language=Python]
numbers = [1, 2, 3, 4, 5]
\end{lstlisting}

and:

\begin{lstlisting}[language=Python]
def change_numbers(nums):

    nums[0] = 7
\end{lstlisting}

Then:

\begin{lstlisting}[language=Python]
change_numbers(numbers)

print(numbers)
\end{lstlisting}

produces:

\begin{lstlisting}
[7, 2, 3, 4, 5]
\end{lstlisting}

Many people incorrectly conclude that Python uses pass-by-reference.

\section{What is Actually Happening}

The important distinction is:

\begin{itemize}
    \item the list itself exists as a mutable object in memory,
    \item the function receives a copy of the reference,
    \item both variables point to the same underlying object.
\end{itemize}

The variables themselves are still separate variables.

\section{Mutable vs Immutable Objects}

This behavior depends heavily on mutability.

\subsection{Immutable Objects}

Examples:

\begin{itemize}
    \item integers
    \item floats
    \item strings
    \item tuples
\end{itemize}

cannot be modified in-place.

\subsection{Mutable Objects}

Examples:

\begin{itemize}
    \item lists
    \item dictionaries
    \item sets
\end{itemize}

can be modified in-place.

\section{Key Takeaways}

\begin{itemize}
    \item Functions help organize and reuse code.

    \item Arguments pass information into functions.

    \item Return values pass information back out.

    \item Python primarily behaves like pass-by-value.

    \item Mutable objects can create reference-sharing behavior.

    \item Scope determines where variables exist and have meaning.

    \item Python uses dynamic typing.
\end{itemize}

\end{document}
